\section{Zadanie 10}
\autor{Kacper Orszulak}
Niech X będzie zmienną losową z rozkładu jednostajnego na przedziale \([0, 1]\). Wykaż, że zmienna losowa \(-ln\;X\) ma rozkład wykładniczy z parametrem 1.

\subsection*{Rozwiązanie}
Niech \(Y := -ln \; X \iff X = e^{-Y}\). Wtedy
\[
	F_Y(y) = \mathcal{P}(Y \le y) = \mathcal{P}(-ln\;X \le y) = \mathcal{P}(X \ge e^{-y}) = 1 - \mathcal{P}(X \le e^{-y}) = 1 - F_X(e^{-y}) = 1 - e^{-y}.
\]

Dla przypomnienia: \(Z \sim Exp(\lambda)\) wtw. \(F_Z(z) = \begin{cases} 1-e^{-\lambda z},\;x \ge 0 \\ 0, \text{ wpp.}\end{cases}\)

Skoro \(img\;X = [0, 1]\), to \(img\;Y = [0, +\infty)\), a zatem rzeczywiście \(Y \sim Exp(1)\). \(\hfill \square\)
